\begin{table}[htbp]
\centering
\caption{Within-Band AUC: FL Screen Performance by Participation Volume}
\label{tab:within_band_auc}
\begin{adjustbox}{max width=\textwidth}
\begin{threeparttable}
\small
\begin{tabular}{clcccc}
\toprule
Decile & Tenders range & Firms & CADE-linked & AUC \\
\midrule
1 & 1--1 & 1,684 & 14 & 0.500 \\
2 & 1--2 & 1,684 & 10 & 0.613 \\
3 & 2--3 & 1,684 & 19 & 0.551 \\
4 & 3--5 & 1,685 & 14 & 0.563 \\
5 & 5--8 & 1,684 & 18 & 0.467 \\
6 & 8--13 & 1,684 & 20 & 0.510 \\
7 & 13--19 & 1,685 & 38 & 0.493 \\
8 & 19--30 & 1,684 & 40 & 0.559 \\
9 & 30--57 & 1,684 & 55 & 0.500 \\
10 & 58--1,126 & 1,685 & 147 & 0.633 \\
\midrule
\multicolumn{4}{l}{Overall AUC (unconditional)} & 0.738 \\
\multicolumn{4}{l}{Pooled within-band AUC} & 0.539 \\
\multicolumn{4}{l}{Residualized AUC (demeaned)} & 0.505 \\
\bottomrule
\end{tabular}
\begin{tablenotes}
\small
\item \textit{Notes:} AUC computed using participation count as the
screening score and CADE co-participation as the binary outcome,
within each decile of the always-loser participation distribution.
Within-band AUC measures discrimination after controlling for volume.
Residualized AUC uses participation count demeaned within deciles.
Total always-losers: 16,843 (CADE-linked: 375).
The ROC analysis in Table~\ref{tab:roc_detection} varies the IQR
\emph{multiplier} rather than the raw count, producing a higher
overall AUC (0.94) because threshold variation spans a wider range
of the participation distribution.
\end{tablenotes}
\end{threeparttable}
\end{adjustbox}
\end{table}
